\PassOptionsToPackage{pdfpagelabels=false}{hyperref}
\documentclass[10pt,a4paper]{ctexart}

\usepackage[left=1.32cm,right=1.32cm,top=1.05cm,bottom=0.08cm]{geometry}
\usepackage{array}
\usepackage{enumitem}
\usepackage{titlesec}
\usepackage{fontawesome5}
\usepackage{xcolor}
\usepackage{hyperref}
\usepackage{bookmark}
\usepackage{eso-pic}

\definecolor{SectionBlue}{HTML}{3E68C6}
\definecolor{MutedGray}{HTML}{5B6472}
\definecolor{KeywordGray}{HTML}{1F3A52}
\definecolor{KeywordGrayZh}{HTML}{1F3A52}

% 页面设置
\pagestyle{empty}
\setlength{\parindent}{0pt}
\setlist[itemize]{leftmargin=1.15em,itemsep=0.10em,topsep=0.08em,parsep=0.01em,partopsep=0em}
\linespread{1.162}
\setmainfont{Times New Roman}
% Let ctex choose the platform-appropriate CJK fonts.

% 超链接样式
\hypersetup{
  colorlinks=true,
  urlcolor=black,
  linkcolor=black,
  pdfnewwindow=true
}

% 分节样式：标题 + 下划线
\titleformat{\section}{\heiti\zihao{3}\color{SectionBlue}}{}{0pt}{}[\vspace{0.06em}{\color{SectionBlue!40}\titlerule}]
\titlespacing*{\section}{0pt}{0.42em}{0.18em}

% 条目：左侧主标题 + 右侧角色/学位；下一行左侧描述 + 右侧时间
\newcommand{\entry}[4]{%
  \if\relax\detokenize{#3}\relax
    \textbf{{\fontsize{11.2pt}{12.2pt}\selectfont #1}}\hfill #2\enspace #4\par
  \else
    \textbf{{\fontsize{11.2pt}{12.2pt}\selectfont #1}}\hfill #2\\
    #3\hfill #4\par
  \fi
}

% 项目标题行：左侧项目名 + 中间标签 + 右侧代码/演示链接
\newcommand{\projentry}[4]{%
  \if\relax\detokenize{#2}\relax
    \makebox[0.72\textwidth][l]{\textbf{{\fontsize{11.2pt}{12.2pt}\selectfont #1}}}%
    \makebox[0.28\textwidth][r]{\codelink{#3}\if\relax\detokenize{#4}\relax\else\enspace\demolink{#4}\fi}\par
  \else
    \makebox[0.52\textwidth][l]{\textbf{{\fontsize{11.2pt}{12.2pt}\selectfont #1}}}%
    \makebox[0.20\textwidth][c]{\textcolor{MutedGray}{\fontsize{11.1pt}{11.9pt}\selectfont\textbf{#2}}}%
    \makebox[0.28\textwidth][r]{\codelink{#3}\if\relax\detokenize{#4}\relax\else\enspace\demolink{#4}\fi}\par
  \fi
}

% 仅将 GitHub 链接显示为蓝色，其他链接保持黑色
\newcommand{\ghlink}[2]{\href[pdfnewwindow=true]{#1}{\textcolor{blue}{#2}}}
\newcommand{\ghstats}[2]{\textcolor{MutedGray}{\small\faStar\ #1\enspace|\enspace\faHistory\ #2}}
\newcommand{\codelink}[1]{\href[pdfnewwindow=true]{#1}{\textcolor{blue}{\faCode\ {\fontsize{10.8pt}{11.2pt}\selectfont\textit{code}}}}}
\newcommand{\demolink}[1]{\href[pdfnewwindow=true]{#1}{\textcolor{blue}{\faVideo\ {\fontsize{10.8pt}{11.2pt}\selectfont\textit{demo}}}}}
\providecommand{\faExternalLinkAlt}{\faIcon{external-link-alt}}
\newcommand{\paperlink}[1]{\nobreak\hspace{0.18em}\href[pdfnewwindow=true]{#1}{\textcolor{MutedGray}{\fontsize{8.8pt}{9.2pt}\selectfont\faExternalLinkAlt}}}
\newcommand{\projkeyword}[1]{\textcolor{KeywordGray}{#1}}
\newcommand{\projkeywordzh}[1]{\textcolor{KeywordGrayZh}{#1}}
\newcommand{\projmetric}[1]{\textcolor{SectionBlue}{#1}}

% 教育条目：学校左对齐，专业/学历状态居中，时间右对齐
\newcommand{\eduentry}[4]{%
  \makebox[0.40\textwidth][l]{\textbf{#1}}%
  \makebox[0.25\textwidth][c]{#2}%
  \makebox[0.15\textwidth][c]{#3}%
  \makebox[0.20\textwidth][r]{#4}\par
}

\begin{document}
\AddToShipoutPictureFG*{%
  \AtPageLowerLeft{%
    \hspace*{1.32cm}\raisebox{0.25cm}[0pt][0pt]{%
      {\fontsize{8pt}{11pt}\selectfont\textcolor{MutedGray}{\begin{tabular}{@{}l@{}}
      \end{tabular}}}%
    }%
  }%
}

% ===== 顶部信息 =====
\vspace*{-1.52cm}
\begin{center}
  {\zihao{2}\bfseries 张三}\\[0.4em]
  {\fontsize{11.8pt}{12.3pt}\selectfont
  \faGlobe\ \ghlink{https://ascendho.netlify.app/}{{\fontsize{10.2pt}{10.7pt}\selectfont 个人主页}}
  \enspace \enspace
  \faGithub\ \ghlink{https://github.com/ascendho}{ascendho}
  \enspace \ghstats{0}{0}
  \enspace \enspace
  \faEnvelope\ \href{mailto:ascendho@example.com}{ascendho@example.com}
  \enspace \enspace
  \faPhone\ +86\ 138\ 0000\ 0000}
\end{center}
\vspace*{-0.48em}

% ===== 教育经历 =====
\section{教育背景}
\eduentry{某某大学·计算机学院}{软件工程}{硕士在读}{2024.09 -- 2027.06}
\eduentry{某某大学·计算机学院}{计算机科学与技术}{本科}{2020.09 -- 2024.06}

% ===== 专业技能 =====
\section{专业技能}
\begin{itemize}
  \item 理解向量嵌入的底层原理，具备使用 Qdrant、Redis 等主流向量数据库构建高效检索系统的实操经验
  \item 深入理解 LLM Agent 架构与 RAG 检索技术，具备使用 LangChain 与 LangGraph 构建复杂的 Agent 工作流的经验
  \item 熟悉模型对齐微调流水线，了解 SFT 与 LoRA 技术，具备基于 GRPO 强化学习算法优化模型推理能力的实战经验
  \item \textbf{英语能力：}CET-4 \textbf{512} 分，CET-6 \textbf{540} 分，2024 年考研英语二 \textbf{85} 分，具备适应英文工作环境的能力
\end{itemize}

% ===== 项目经历 =====
\section{项目经历}
\projentry{E-Snap 两级缓存与语义增强的智能电商客服}{}{https://github.com/ascendho/E-Snap}{https://ascendho.netlify.app/projects/e-snap/demo}
\begin{itemize}
  \item \textbf{技术栈：} Python、FastAPI、LangGraph、LangChain、RedisVL / Redis Stack
  \item \textbf{项目介绍：}面向电商客服场景的语义缓存增强系统，重点解决高频问题重复检索、动态问题误答风险等问题。系统启动时先对高频 FAQ 进行预加载缓存，知识检索侧采用向量召回与关键词检索结合的\projkeywordzh{混合检索机制}。请求进入后，先对时间、库存、具体商品型号等动态查询进行前置拦截，避免不适合静态复用的问题进入缓存。随后进入\projkeywordzh{两级缓存}：L1 通过归一化、编辑距离和子问题命中等规则优先复用历史答案，L2 通过\projkeywordzh{语义缓存}召回相似问答；对于存在候选但不能直接复用的问题，再由 \projkeyword{LLM Validator} 判断是直接复用、补充缺失信息，还是回退到完整 RAG。只有在缓存无法覆盖或复用收益不足时，系统才进入 \projkeyword{ReAct} 式检索完成知识查询和答案生成，并将结果按规则写回缓存。
  \item \textbf{实验评估：}在 30 条离线评测请求上，系统前置拦截 4 条动态类查询，缓存覆盖率达到 42.31\%，其中直接缓存复用率为 30.77\%，使端到端延迟降低了 \projmetric{26.64\%}，系统吞吐量提升至原来的 \projmetric{1.36} 倍；成本方面，相对纯 RAG 基线，平均每次请求减少 1.14 次 LLM 调用，累计净节省约 25 次调用，费用总体节省 \projmetric{32.36\%}。
\end{itemize}

\projentry{VisionDoc 视觉检索驱动的多模态页证智答系统}{}{https://github.com/ascendho/VisionDoc}{https://ascendho.netlify.app/projects/visiondoc/demo}
\begin{itemize}
  \item \textbf{技术栈：}Python、FastAPI、ColPali、MUVERA、Qdrant、pdf2image / Poppler、Pillow、Docker
  \item \textbf{项目介绍：}聚焦传统 OCR 流程难以稳定处理的图文混排、复杂版式、表格截图与页面结构信息丢失等问题，构建多模态页级证据检索与问答系统。系统在上传阶段先将 PDF、PPTX 与图片统一转换为页面图像，随后使用 ColPali 为每一页生成多向量视觉特征，再使用 MUVERA 生成压缩召回向量，并将两类向量一并写入 Qdrant 建立视觉索引。查询阶段，后端先进行无关问题过滤，再进行\projkeywordzh{复合问题拆分}与\projkeywordzh{多轮对话上下文改写}，接着通过 \projkeyword{MUVERA Prefetch + ColPali MaxSim} 精排执行两阶段视觉检索。若高置信证据量不足，则补入 near-threshold 页面并在必要时回退到 top-1 候选页，最后通过豆包多模态模型生成答案，返回带文档名、页码和页面截图溯源的结果。
  \item \textbf{实验评估：}在包含 3 份真实文档的 20 题测评中，双路检索机制展现出显著的设计优势：相比单路（muvera\_only）基线，该架构在平均检索延迟仅微增约 28ms（至 444.3ms）的情况下，将 Hit Rate@3 由 25\% 大幅跃升至 \projmetric{90\%}，MRR 由 0.25 提升至 \projmetric{0.7667}，并实现了 \projmetric{1.00} 的 Faithfulness，表明系统能够以极低的时间开销换取极其稳健的定位精度。
\end{itemize}

\projentry{基于 SFT+GRPO 的 Wordle 反馈驱动两阶段微调}{}{https://github.com/ascendho/Wordle}{}
\begin{itemize}
  \item \textbf{技术栈：}Python、Predibase、LoRA、SFT、GRPO
  \item \textbf{项目介绍：}针对通用模型在结构化输出不稳定、历史反馈利用不足和多轮搜索策略薄弱的问题，围绕 Wordle（一种猜词游戏）搭建基于 \textit{Qwen 2.5 7B Instruct} 与 Predibase 的两阶段微调流水线。首先通过 SFT 对齐 think 和 guess 双段结构化输出与基础规则推理能力，再通过 GRPO 引入\projkeywordzh{格式校验}、\projkeywordzh{历史反馈一致性}、\projkeywordzh{信息增益}三类奖励函数，强化模型对反馈的利用和候选空间压缩能力，并构建真实对局 Benchmark 用于评估模型表现。
  \item \textbf{实验评估：}在 10 局随机基准测试、单局最多 6 次猜测的设定下，相较基础 \textit{Qwen 2.5 7B Instruct}（0/10 通关）和单 GRPO 方案（3/10 通关，4.0/6 步），SFT + GRPO 联合训练方案取得 \projmetric{7/10} 通关、获胜局平均 \projmetric{4.0/6} 步的优异成绩，表现\projkeywordzh{超过}闭源大模型\textit{GPT-4o-mini}（1/10 通关，4.0/6 步），并\projkeywordzh{逼近} \textit{Claude 3.5 Sonnet}（8/10 通关，3.9/6 步），验证了小参数模型在垂直领域推理任务上的强化学习增益。
\end{itemize}

% ===== 科研成果 =====
\section{科研成果}
\begin{itemize}
  \item \textbf{SCI 二区在投，学生一作：}物联网异常检测框架研究
  \item \textbf{IEEE 某国际会议，第二作者：}跨设备数据安全验证方法研究
  \item \textbf{一项发明专利}已受理；\textbf{一项软件著作权}已登记
\end{itemize}

% ===== 证书 =====
\section{奖项与证书}
\begin{itemize}
  \item \textbf{MOOC 认证：}自主学习大模型应用开发最新技术，并在 Coursera、Udemy、Harvard 等平台取得\ghlink{https://ascendho.netlify.app/\#certificates}{多项课程证书}
  \item \textbf{竞赛：}某省级创新大赛（2025）\textbf{铜奖}
  \item \textbf{奖学金：}2024、2025学年研究生三等学业奖学金；本科期间获院级二等、三等奖学金各一次
\end{itemize}

\end{document}
